\begin{center}
  \hyperlink{https://www.kaggle.com/datasets/cnic92/200-financial-indicators-of-us-stocks-20142018}{\textit{200+ Financial Indicators of US stocks (2014-2018)}}
\end{center}

The dataset consists of five annual CSV files spanning 2014 through 2018, each containing approximately 4000 rows and 223 financial indicator columns. Rows correspond to individual company-year observations identified by a ticker symbol and a sector label. The dataset was assembled from the Financial Modeling Prep API and pandas-datareader and is available on Kaggle under the title 200+ Financial Indicators of US Stocks (2014-2018).

Across the five files, 4980 distinct ticker symbols appear, though not every ticker is present in every year — the year-over-year overlap is roughly $74.82\%$, reflecting delistings, IPOs, and coverage gaps in the source API. Financial Services is the most heavily represented sector in every year, accounting for approximately $21.38\%$ of unique tickers annually, followed by Healthcare and Technology. The remaining sectors — Industrials, Consumer Cyclical, Basic Materials, Real Estate, Utilities, Energy, Consumer Defensive, and Communication Services — each contribute a smaller and comparatively stable share.

Missing values are pervasive. There are roughly 5 million cells; $16.61\%$ are \texttt{NaN}. Columns representing less common indicators (e.g., net debt to EBITDA, free cash flow yield) have substantially higher \texttt{NaN} rates than core metrics like revenue or EPS. The author of the dataset notes the presence of outliers likely caused by data-entry errors; inspection confirms that several indicators contain values orders of magnitude larger than the interquartile range, which motivated the winsorization step often used in the analytics agent's preprocessing pipeline.

These characteristics shaped several architectural decisions. The \texttt{NaN} density made it impractical to require all 223 columns to be populated for a row to be valid; instead, the analytics agent selects and cleans only the columns relevant to each query. The outlier distribution motivated the use of median rather than mean aggregation in sector-level comparisons, and explicit winsorization before z-score normalization in composite scoring tasks.
